\documentclass[10pt,a4paper]{amsart}
\usepackage[margin=19mm]{geometry}
\usepackage{amsmath,amssymb,mathtools}
\usepackage[scr=rsfs,cal=euler]{mathalfa}
\usepackage{xfrac}
\usepackage[unicode,hidelinks,bookmarks=false]{hyperref}
\newcommand{\ie}{i.e.\ }
\newcommand{\cf}{cf.\ }
\newcommand{\resp}{respectively\ }
\newcommand{\Subsection}{\subsection}
\providecommand{\nobreakdash}{\nobreak\hbox{-}}
\setlength{\emergencystretch}{2em}
\multlinegap0pt
\usepackage{csquotes}
\usepackage{amssymb}
\usepackage{mathtools}
\usepackage{esint}
\usepackage{float}
\usepackage{tikz}
\usepackage{caption}
\newtheorem{lemma}{Lemma}
\newtheorem{theorem}{Theorem}
\title{Morse index of min--max stationary integral varifolds}
\author{Mitchell Gaudet, Talant Talipov}
\date{Source: arXiv:2601.02860v1 (CC BY 4.0)}
\begin{document}
\maketitle

\noindent\emph{Source and licence.} Morse index of min--max stationary integral varifolds. Version arXiv:2601.02860v1 (CC BY 4.0), distributed by arXiv under the Creative Commons Attribution 4.0 International licence, http://creativecommons.org/licenses/by/4.0/, as recorded in the arXiv licence metadata for this exact version and shown on the abstract page https://arxiv.org/abs/2601.02860v1. Full licence notice: \texttt{licenses/arxiv-2601.02860-CC-BY-4.0.txt}.\par\smallskip

\noindent\textit{Editorial scope.} Counted unit: the article's theorem Thm:deform (source0.tex lines 864--875). The article's complete original proof of it is delivered with it (source0.tex lines 876--1280, one \texttt{proof} environment of 405 lines). No mathematics is written, repaired or re-derived: the statement and the proof are the article's own lines, in the article's own notation, and no retained line is substituted or re-flowed.  Retained uncounted as the source-local context the proof needs: lines 457--467; lines 536--545; lines 572--578; lines 635--646; lines 805--828.  One declared adaptation: exactly 2 ancillary strings inside the retained spans are omitted (image-inclusion commands and, where the source repeats a label, the repeated label definitions) and no other line differs from the source; the figure environments, with their placement, captions and labels, are kept, and the package records the omitted strings in fidelity receipts bound to the affected bodies.  No retained line contains a citation, so the wrapper carries no bibliography and no bibliography entry.  The retained lines contain the article's own diagram code, which the wrapper typesets with the standard tikz packages; no image file is delivered and the package declares no asset dependency.  Editorial formatting only, in addition: an AMS \texttt{amsart} preamble that restates the article's own declarations and macros used by the retained lines, namely the environments \texttt{lemma}, \texttt{theorem} on separate counters, together with the standard packages those lines need.  The retained environments print in this document as Lemma 1, Theorem 1 in order of appearance, whereas the article prints the same items with the numbering of its own full document.  The complete native source of the article as distributed in the arXiv e-print for this exact version is bundled unchanged as \texttt{sources/192188/source0.tex}.
\begin{lemma}
\label{lem:grad}
    Given any fixed $V \in \mathbf{\overline{B}}^{\textbf{F}}_{2\varepsilon}(S)$, if an integral curve $x: [0, L] \rightarrow \overline{B}^k$ disjoint from $m(V)$ is defined by
    \begin{equation*}
        \frac{d}{dt} x(t) = - f(t) \nabla A^V\bigl(x(t)\bigr),
    \end{equation*}
    where $f(t)$ is a positive continuous function, then we have
    \begin{equation*}
        A^V\bigl(x(L)\bigr) - A^V\bigl(x(0)\bigr) \leq - \frac{c_0}{2} |x(L) - x(0)|^2.
    \end{equation*}
\end{lemma}

\begin{lemma} \label{lem:dist_F}
    There exists an increasing function
    \begin{equation*}
        h_S : \mathbb{R}^+ \rightarrow \mathbb{R}^+,
    \end{equation*}
    such that for any $V \in \overline{\mathbf{B}}^{\mathbf{F}}_{2\varepsilon}(S)$, and $v, w \in \overline{B}^k$,
    \begin{equation*}
        |v - w| \geq h_S \Bigl(\mathbf{F}\bigl((F_v)_\#(V), (F_w)_\#(V)\bigr)\Bigr)
    \end{equation*}
\end{lemma}

\begin{lemma}\label{lem:union_compact}
    The set \(\mathcal{U}(\omega^d_p)\) is a countable union of compact sets, i.e.
    \[
        \mathcal{U}(\omega^d_p)=\bigcup_{i=1}^\infty\mathcal{U}_i(\omega^d_p),
    \]
    where each \(\mathcal{U}_i(\omega^d_p)\) is compact.
\end{lemma}

\begin{lemma}\label{lem:gradflow}
Let \(\lambda>0\) and \(K\subset\mathbb N^+\) be finite, and take \(\eta>0\) sufficiently small. Then there exist constants \(c(\lambda,K)>0\) and \(T>0\) such that the following holds: for any \(V\in\overline{\mathbf{B}}^{\mathbf F}_{\lambda,K}\) and any \(v\in\overline{B}^{p(n-d)+1}\) with \(|v-m_{\bar{k}(K)}(V)|\ge \eta\) and
\[
(F_v^{\bar{k}(K)})_\#(V)\in \overline{\mathbf{B}}^{\mathbf F}_{\lambda+0.5,K},
\]
we have
\[
A^V_{\bar{k}(K)}\bigl(\phi^V_{\lambda,K}(v,T)\bigr)
<
A^V_{\bar{k}(K)}\bigl(\phi^V_{\lambda,K}(v,0)\bigr)-c(\lambda,K).
\]
\end{lemma}

\begin{lemma}\label{lem:homotopy}
Let \(\lambda>0\), and let \(K\subset \mathbb{N}^+\) be a nonempty finite set with \(\mathbf{B}^\mathbf{F}_{\lambda,K}\neq\emptyset\). Let \(c(\lambda,K)\) be the constant from Lemma~\ref{lem:gradflow}. For any \(\delta\in(0,c(\lambda,K)/4)\), any \(l\in\{0,1,\dots,p(n-d)\}\), and any continuous map
\[
\Theta:X^{l}\to \mathbf{B}^\mathbf{F}_{\lambda,K},
\]
where \(X^l\) is an \(l\)-dimensional finite simplicial complex, there exists a homotopy
\[
H^{\Theta,\delta}_{\lambda,K}:X^l\times[0,1]\to \mathbf{B}^\mathbf{F}_{\lambda+1,K}
\]
such that:
\begin{enumerate}
    \item \(H^{\Theta,\delta}_{\lambda,K}(\cdot,0)=\Theta(\cdot)\).
    \item For all \(x\in X^l\), \(\ \|H^{\Theta,\delta}_{\lambda,K}(x,1)\|(M)\le \|\Theta(x)\|(M)-\tilde c(\lambda,K)\).
    \item For all \(x\in X^l\) and \(t\in[0,1]\), \(\ \|H^{\Theta,\delta}_{\lambda,K}(x,t)\|(M)-\|\Theta(x)\|(M)\le \delta\).
\end{enumerate}
Here \(\tilde c(\lambda,K):=c(\lambda,K)/2\). More precisely:
\begin{itemize}
    \item For \(t\in[0,1/2]\), \(\ \mathbf{F}\bigl(H^{\Theta,\delta}_{\lambda,K}(x,t),\Theta(x)\bigr)\le \delta\).
    \item For \(t\in[1/2,1]\) there exists \(v(x)\in \overline{B}^{p(n-d)+1}\) such that
    \[
        H^{\Theta,\delta}_{\lambda,K}(x,t)=\phi^{\Theta(x)}_{\lambda,K}\bigl(v(x),(2t-1)T\bigr).
    \]
\end{itemize}
\end{lemma}

\begin{theorem}[Deformation Theorem]\label{Thm:deform}
Let \(\{\Phi_i\}_{i\in\mathbb N}\) be a min--max sequence for \(\omega_p^d\), and let \(X_i=\operatorname{dmn}(\Phi_i)\) be its domain with \(\dim X_i\le p(n-d)\). After passing to a subsequence, there exists a sequence \(\{\Psi_i\}_{i\in\mathbb N}\) such that:
\begin{enumerate}
    \item \(\Psi_i\) is homotopic to \(\Phi_i\) in the \(\mathbf F\)-topology;
    \item \(\mathbf L\bigl(\{\Psi_i\}_{i\in\mathbb N}\bigr)=\omega^d_p\);
    \item There exists a function \(\bar\varepsilon:\mathcal U(\omega^d_p)\to(0,\infty)\) with the following property: for every \(S\in\mathcal U(\omega^d_p)\) there exists \(i_0\in\mathbb N\) such that, for all \(i\ge i_0\),
    \[
        \bigl|\Psi_i\bigr|(X_i)\ \cap\ \mathbf B^{\mathbf F}_{\bar\varepsilon(S)}(S)\;=\;\emptyset .
    \]
    In particular, \(\mathbf C(\{\Psi_i\})\cap \mathcal{APR}_{p,d} \subset \mathcal S(\omega^d_p)\).
\end{enumerate}
\end{theorem}

\begin{proof}
The proof consists of 4 steps.

\textbf{Step 1.} We construct a (finite or countable) sequence of \(\bigl(p(n-d)+1\bigr)\)-unstable stationary integral varifolds
\(\{S_k\}\subset\mathcal{U}(\omega^d_p)\), together with associated quadruples
\(\{(\varepsilon_k,c_{0,k},\{F^k_v\},m_k)\}\), so that the following hold:

\begin{itemize}
    \item For \(k_1\neq k_2\) we have \(\mathbf B^\mathbf F_{\varepsilon_{k_1}}(S_{k_1})\not\subset \mathbf B^\mathbf F_{\varepsilon_{k_2}}(S_{k_2})\).
    \item \(\mathcal{U}(\omega^d_p)\subset\bigcup_{k=1}^\infty \mathbf B^\mathbf F_{\varepsilon_k}(S_k)\), and moreover, for every \(m\in\mathbb N\) there exists \(k_m\in\mathbb N\) with
    \[
        \mathcal{U}_m(\omega^d_p)\subset\bigcup_{i=1}^{k_m}\mathbf B^\mathbf F_{\varepsilon_i}(S_i).
    \]
    \item For each fixed \(k\), the set
    \[
        \bar K(S_k):=\{k':\ \mathbf B^\mathbf F_{2\varepsilon_{k'}}(S_{k'})\cap\mathbf B^\mathbf F_{2\varepsilon_k}(S_k)\neq\emptyset\}
    \]
    is finite.
\end{itemize}

Indeed, for every \(S\in\mathcal{U}(\omega^d_p)\) choose a quadruple \((\varepsilon_S,c_{0,S},\{F^S_v\},m_S)\). The family \(\{\mathbf B^\mathbf F_{\varepsilon_S}(S)\}_{S\in\mathcal{U}(\omega^d_p)}\) covers \(\mathcal{U}(\omega^d_p)\). Note that one may always replace \(\varepsilon_S\) by a smaller positive number and keep \(S\) \(\bigl(p(n-d)+1\bigr)\)-unstable with the new quadruple.

If \(\mathcal{U}(\omega^d_p)\) is compact, we extract a finite subcovering \(\{\mathbf B^\mathbf F_{\varepsilon_k}(S_k)\}_{k=1}^{k_1}\) and take the corresponding quadruples. Then all the required properties hold.

Otherwise, by Lemma~\ref{lem:union_compact} we have a decomposition
\[
    \mathcal{U}(\omega^d_p)=\bigcup_{m\in\mathbb N}\mathcal{U}_m(\omega^d_p),
\qquad
\mathcal{U}_m(\omega^d_p):=\{S\in\mathcal{U}(\omega^d_p):\ \mathbf F(S,\mathcal{S}(\omega^d_p))\ge 1/m\},
\]
with each \(\mathcal{U}_m(\omega^d_p)\) compact. We construct the sequence \(\{S_k\}\) inductively in \(m\).

For \(m=1\): for every \(S\in\mathcal{U}_1(\omega^d_p)\) choose \(\varepsilon_S\) so that \(\varepsilon_S<1/4\). Extract a finite subcovering \(\{\mathbf B^\mathbf F_{\varepsilon_k}(S_k)\}_{k=1}^{k_1}\) of \(\mathcal{U}_1(\omega^d_p)\) with minimal \(k_1\). If \(V\in\mathbf B^\mathbf F_{2\varepsilon_k}(S_k)\) for some \(k\le k_1\), then
\[
    \mathbf F(V,\mathcal{S}(\omega^d_p)) \ge \mathbf F(S_k,\mathcal{S}(\omega^d_p)) - \mathbf F(S_k,V)
    \ge 1 - 2\varepsilon_k \ge 1/2.
\]
Hence
\[
    \bigcup_{k=1}^{k_1}\mathbf B^\mathbf F_{2\varepsilon_k}(S_k) \subset \mathcal{U}_2(\omega^d_p).
\]

Inductively, suppose we have chosen \(\{S_k\}_{k=1}^{k_{m-1}}\) covering \(\mathcal{U}_{m-1}(\omega^d_p)\). Set
\[
    \mathcal R_m := \mathcal{U}_m(\omega^d_p)\setminus\bigcup_{k=1}^{k_{m-1}}\mathbf B^\mathbf F_{\varepsilon_k}(S_k),
\]
which is compact. For each \(S\in\mathcal R_m \), choose \(\varepsilon_S<\tfrac{1}{2m(m+1)}\), and then take a finite subcovering \(\{\mathbf B^\mathbf F_{\varepsilon_k}(S_k)\}_{k=k_{m-1}+1}^{k_m}\) with minimal \(k_m\). Then for any \(k\le k_m\) and \(V\in\mathbf B^\mathbf F_{2\varepsilon_k}(S_k)\) we have
\[
    \mathbf{F}\bigl(V, \mathcal{S}(\omega^d_p)\bigr) \geq \mathbf{F}\bigl(S_k, \mathcal{S}(\omega^d_p)\bigr) - \mathbf{F}(S_k, V) \geq \frac{1}{m} - 2\varepsilon_k \geq \frac{1}{m} - \frac{1}{m(m+1)} = \frac{1}{m+1}.
\]
Moreover, if \(k\ge k_{m-1}+1\), then
\begin{align*}
    \mathbf{F}\bigl(V, \mathcal{S}(\omega^d_p)\bigr) &\leq \mathbf{F}\bigl(S_k, \mathcal{S}(\omega^d_p)\bigr) + \mathbf{F}(S_k, V) \\
    &\leq \frac{1}{m-1} + 2\varepsilon_k \leq \frac{1}{m-1} + \frac{1}{m(m+1)} \\
    &< \frac{1}{m-1} + \frac{1}{m+1} = \frac{2m}{m^2 - 1} < \frac{2}{m} < \frac{1}{\lfloor m/2\rfloor}.
\end{align*}
Thus
\begin{equation}\label{eq:subset1}
    \bigcup_{k=1}^{k_m}\mathbf B^\mathbf F_{2\varepsilon_k}(S_k) \subset \mathcal{U}_{m+1}(\omega^d_p),
\end{equation}
and
\begin{equation}\label{eq:subset2}
    \bigcup_{k=k_{m-1}+1}^{k_m}\mathbf B^\mathbf F_{2\varepsilon_k}(S_k)
    \subset \mathcal{U}_{m+1}(\omega^d_p)\setminus\mathcal{U}_{\lfloor m/2\rfloor}(\omega^d_p).
\end{equation}

Continuing this process for all \(m\) yields the desired sequence \(\{S_k\}_{k\in\mathbb N}\). The minimality in each stage guarantees the first bullet. The covering property is built into the construction, giving the second bullet. By~\eqref{eq:subset1} and~\eqref{eq:subset2}, for a fixed \(k\le k_m\), any ball \(\mathbf B^\mathbf F_{2\varepsilon_{k'}}(S_{k'})\) that intersects \(\mathbf B^\mathbf F_{2\varepsilon_k}(S_k)\) must itself be among the finitely many balls produced up to stage \(2m+2\). Hence
\[
    \#\bar K(S_k)\le k_{2m+2}<\infty,
\]
and the third bullet follows.


\textbf{Step 2.} We construct a function \(\eta:\mathcal{U}(\omega^d_p)\to(0,1]\) such that:
\begin{itemize}
    \item The neighborhoods satisfy
    \[
    \begin{aligned}
    \mathcal N_\eta &:= \bigcup_{S\in\mathcal U(\omega^d_p)} \mathbf B^\mathbf{F}_{2\eta(S)}(S)
    \subset \bigcup_{k} \overline{\mathbf B}^\mathbf{F}_{\varepsilon_k}(S_k),\\
    \mathcal N^m_\eta &:= \bigcup_{S\in\mathcal U_m(\omega^d_p)} \mathbf B^\mathbf{F}_{2\eta(S)}(S)
    \subset \bigcup_{k=1}^{k_m} \overline{ \mathbf B}^\mathbf{F}_{\varepsilon_k}(S_k).
    \end{aligned}
    \]
    \item On each \(\mathcal{U}_m(\omega^d_p)\), the function \(\eta\) has a positive lower bound \(\tilde\eta_m>0\).
    \item For any \(\lambda\in\{0,1,\dots,p(n-d)\}\), \(l\in\{0,1,\dots,p(n-d)\}\), \(K\subset\mathbb{N}\), \(\delta\in\bigl(0,c(\lambda,K)/4\bigr)\), and any continuous map \(\Theta:X^l\to \mathcal{N}_\eta\cap \mathbf{B}^\mathbf{F}_{\lambda,K}\) (assume \(\mathcal{N}_\eta\cap \mathbf{B}^\mathbf{F}_{\lambda,K}\neq\emptyset\)), there exists a homotopy
    \[
        H^{\Theta,\delta}_{\lambda,K}:X^l\times[0,1]\to \mathbf{B}^\mathbf{F}_{\lambda+1,K}
    \]
    such that:
    \begin{enumerate}
        \item \(H^{\Theta,\delta}_{\lambda,K}(\cdot,0)=\Theta(\cdot)\).
        \item \(\|H^{\Theta,\delta}_{\lambda,K}(x,t)\|(M)-\|\Theta(x)\|(M)\le \delta\) for all \(t\in[0,1]\).
        \item For \(t\in[0,1/2]\), \(\mathbf{F}\bigl(H^{\Theta,\delta}_{\lambda,K}(x,t),\Theta(x)\bigr)\le \delta\).
        \item For \(t\in[1/2,1]\) there exists \(v(x)\in\overline{B}^{p(n-d)+1}\) such that
        \begin{equation}\label{eq:flow}
            H^{\Theta,\delta}_{\lambda,K}(x,t)=\phi^{\Theta(x)}_{\lambda,K}\bigl(v(x),(2t-1)T\bigr).
        \end{equation}
        \item \(H^{\Theta,\delta}_{\lambda,K}(\cdot,1)\notin \mathcal{N}_\eta\).
    \end{enumerate}
\end{itemize}
        
\textbf{Claim 1.} There exists a \emph{continuous} positive function \(\eta_1:\mathcal{U}(\omega^d_p)\to\mathbb{R}^+\) such that, for any \(m\in\mathbb{N}\) and any \(S\in\mathcal{U}_m(\omega^d_p)\),
\[
    \mathbf{B}^\mathbf{F}_{2\eta_1(S)}(S)\subset \bigcup_{k=1}^{k_m}\mathbf{B}^\mathbf{F}_{\varepsilon_k}(S_k).
\]

\begin{proof}
Define \(\eta_1\) inductively. If \(S\in\mathcal{U}_1(\omega^d_p)\), set
\[
    \eta_1(S):=\max\Bigl\{r>0:\ \mathbf{B}^\mathbf{F}_{2r}(S)\subset \bigcup_{k=1}^{k_1}\overline{\mathbf{B}}^\mathbf{F}_{\varepsilon_k}(S_k)\Bigr\}.
\]
If \(m\ge2\) is the least integer with \(S\in \mathcal{U}_m(\omega^d_p)\setminus \mathcal{U}_{m-1}(\omega^d_p)\), define
\[
    \eta_1(S):=\min\left\{
    \begin{array}{l}
    \displaystyle \max\Bigl\{r>0:\ \mathbf{B}^\mathbf{F}_{2r}(S)\subset \bigcup_{k=1}^{k_m}\overline{\mathbf{B}}^\mathbf{F}_{\varepsilon_k}(S_k)\Bigr\},\\[6pt]
    \displaystyle \min\bigl\{\eta_1(S')+\mathbf{F}(S,S'):\ S'\in \mathcal{U}_{m-1}(\omega^d_p)\bigr\}
    \end{array}\right\}>0.
\]
This \(\eta_1\) is continuous and therefore has a positive lower bound on each compact set \(\mathcal{U}_m(\omega^d_p)\).
\end{proof}

Since $\eta_1$ is continuous, the set $\{t\geq0: \forall\, V\in \overline{\mathbf{B}}^\mathbf{F}_{t}(S),\ \forall\,\lambda\in\{0,1,\dots,p(n-d)\},\bigl|\omega^d_p-\|V\|(M)\bigr|\le \tilde c(\lambda,\tilde K)/3\}$ is closed, and since $\|V\|(M)=\omega^d_p$, by continuity of mass we have the supremum of this set is positive. Therefore, we can define \(\eta(S)\) to be the largest number in \(\bigl(0,\min(\eta_1(S),1)]\bigr)\) such that
\[
    \forall\, V\in \mathbf{B}^\mathbf{F}_{\eta(S)}(S),\ \forall\,\lambda\in\{0,1,\dots,p(n-d)\}:\quad
    \bigl|\omega^d_p-\|V\|(M)\bigr|\le \tilde c(\lambda,\tilde K)/3
\]
for every nonempty subset \(\tilde K\subset \{\,k\mid \mathbf{B}^\mathbf{F}_{\eta_1(S)}(S)\cap \mathbf{B}^\mathbf{F}_{2\varepsilon_k}(S_k)\neq\emptyset\,\}\).

To show that \(\eta\) has a positive lower bound \(\tilde\eta_m\) on each \(\mathcal{U}_m(\omega^d_p)\), observe that for any \(S\in\mathcal{U}_m(\omega^d_p)\) the index set
\[
    \{\,k : \mathbf{B}^\mathbf{F}_{\eta_1(S)}(S)\cap \mathbf{B}^\mathbf{F}_{2\varepsilon_k}(S_k)\neq\emptyset\,\}
    \subset K_m:=\bigcup_{k=1}^{k_m}\bar K(S_k),
\]
since
\[\displaystyle \bigcup_{S\in\mathcal{U}_m(\omega^d_p)}\mathbf{B}^\mathbf{F}_{2\eta_1(S)}(S)\subset \bigcup_{k=1}^{k_m}\overline{\mathbf{B}}^\mathbf{F}_{\varepsilon_k}(S_k).\]
By the third bullet in \textbf{Step 1}, the set \(K_m\) is finite. It follows from the compactness of \(\mathcal{U}_m(\omega^d_p)\) that there exists \(\tilde\eta_m>0\) such that, for any \(S\in\mathcal{U}_m(\omega^d_p)\) and any \(V\in \mathbf{B}^\mathbf{F}_{\tilde\eta_m}(S)\),
\[
    \bigl|\omega^d_p-\|V\|(M)\bigr|\le \min_{\substack{K\subset K_m\\ \lambda\in\{0,1,\dots,p(n-d)\}}}\frac{\tilde c(\lambda,K)}{3}.
\]
Hence \(\eta(S)\ge \tilde\eta_m>0\) on \(\mathcal{U}_m(\omega^d_p)\).

For the third bullet, apply Lemma~\ref{lem:homotopy} to construct the homotopy \(H^{\Theta,\delta}_{\lambda,K}\). Properties (1)–(4) follow directly. To verify (5), suppose, toward a contradiction, that there exist \(x\in X^l\) and \(S\in\mathcal{U}(\omega^d_p)\) with
\[
H^{\Theta,\delta}_{\lambda,K}(x,1)\in \mathbf{B}^\mathbf{F}_{\eta(S)}(S)\cap \mathbf{B}^\mathbf{F}_{\lambda+1,K}.
\]
By the definition of \(\eta\),
\[
\omega^d_p-\|H^{\Theta,\delta}_{\lambda,K}(x,1)\|(M)<\tilde c(\lambda,K)/2.
\]
Since \(H^{\Theta,\delta}_{\lambda,K}(x,0)=\Theta(x)\in \mathbf{B}^\mathbf{F}_{\eta(S)}(S)\cap \mathbf{B}^\mathbf{F}_{\lambda+1,K}\), we also have
\[
\|\Theta(x)\|(M)-\omega^d_p<\tilde c(\lambda,K)/2.
\]
Therefore,
\[
\|\Theta(x)\|(M)-\|H^{\Theta,\delta}_{\lambda,K}(x,1)\|(M)<\tilde c(\lambda,K),
\]
contradicting item (2) of Lemma~\ref{lem:homotopy}. This proves (5).

\textbf{Step 3.} We show that the homotopy map defined in \textbf{Step 2} does not push a varifold too close to $\mathcal{U}(\omega^d_p)$, provided it is initially far from $\mathcal{U}(\omega^d_p)$.

More precisely, we claim that for any $S\in\mathcal{U}(\omega^d_p)$ there exist sequences $\{\varepsilon_q(S)\}_{q=1}^{p(n-d)}\subset\mathbb{R}^+$ and $\{a_q(S)\}_{q=1}^{p(n-d)+1}$ with $a_1(S)=2$ and $a_q(S)\ge 2^q$ such that the following holds. In the third bullet of \textbf{Step 2}, for any $q\in\{1,\dots,p(n-d)\}$, any $S\in\mathcal{U}(\omega^d_p)$ with $\mathbf{B}^{\mathbf F}_{\lambda+1,K}\cap \mathbf{B}^{\mathbf F}_{\eta(S)}(S)\neq\emptyset$, and any $x\in X^l$ with $\Theta(x)\notin \mathbf{B}^{\mathbf F}_{\eta(S)/a_q(S)}(S)$ and $\|\Theta(x)\|(M)-\omega^d_p\le \varepsilon_q(S)$, we have for all $t\ge0$ and all $\delta\in\bigl(0,\min\bigl(\eta(S)/(4a_q(S)),\varepsilon_q(S)\bigr)\bigr)$,
\[
    H^{\Theta,\delta}_{\lambda,K}(x,t)\notin \mathbf{B}^{\mathbf F}_{\eta(S)/a_{q+1}(S)}(S).
\]
Moreover, on each $\mathcal{U}_m(\omega^d_p)$ the quantities $\varepsilon_q(S)$ admit a positive lower bound $\tilde\varepsilon^m_q$, and $a_q(S)$ admit a uniform upper bound $\tilde a^m_q$.

We construct $a_q(S)$ and $\varepsilon_q(S)$ inductively. Set $a_1(S)=2$.

Fix $S\in\mathcal{U}(\omega^d_p)$ and suppose $a_q(S)$ is defined (and if $q>1$, also $\varepsilon_{q-1}(S)$). For each \(k\) with \(\overline{\mathbf{B}}^{\mathbf F}_{\lambda+1,\{k\}}\cap \mathbf{B}^{\mathbf F}_{\eta(S)}(S)\neq\emptyset\) and for each \(V\in\overline{\mathbf{B}}^{\mathbf F}_{\lambda+1,\{k\}}\), set
\[
\begin{aligned}
A_{q,S,k}(V)&:=\bigl((F^k_{\cdot})_{\#}V\bigr)^{-1}\partial\mathbf{B}^{\mathbf F}_{\,3\eta(S)/(4a_q(S))}(S),\\
B_{q,S,k}(V)&:=\bigl((F^k_{\cdot})_{\#}V\bigr)^{-1}\overline{\mathbf{B}}^{\mathbf F}_{\,\eta(S)/(2a_q(S))}(S).
\end{aligned}
\]
Define
\[
d_{q,S,k}(V):=\mathrm{dist}\bigl(A_{q,S,k}(V),\,B_{q,S,k}(V)\bigr),
\]
with the convention that \(\mathrm{dist}(\emptyset,\cdot)=\mathrm{dist}(\cdot,\emptyset)=+\infty\). By Lemma~\ref{lem:dist_F},
\[
    d_{q,S,k}(V)\ \ge\ h_{S_k}\!\left(\frac{\eta(S)}{4a_q(S)}\right).
\]
Choose $m \in \mathbb{N}$ such that $S \in \mathcal{U}_m(\omega^d_p)$. The index set
\[
    \{\,k : \overline{\mathbf{B}}^\mathbf{F}_{\lambda+1,\{k\}}\cap \mathbf{B}^\mathbf{F}_{\eta(S)}(S)\neq\emptyset\,\}
    \subset K_m=\bigcup_{k=1}^{k_m}\bar K(S_k),
\]
since
\[\displaystyle \bigcup_{S\in\mathcal{U}_m(\omega^d_p)}\mathbf{B}^\mathbf{F}_{2\eta_1(S)}(S)\subset \bigcup_{k=1}^{k_m}\mathbf{B}^\mathbf{F}_{\varepsilon_k}(S_k).\]
By the third bullet in \textbf{Step 1}, the set \(K_m\) is finite. Hence we can set
\begin{equation}\label{eq:def-eps}
    \varepsilon_q(S):=\min_{\substack{k:\\ \overline{\mathbf{B}}^{\mathbf F}_{\lambda+1,\{k\}}\cap \mathbf{B}^{\mathbf F}_{\eta(S)}(S)\neq\emptyset}}
    \frac{c_{0,k}}{16}\left[h_{S_k}\!\left(\frac{\eta(S)}{4a_q(S)}\right)\right]^2.
\end{equation}

\textbf{Claim 2.} There exists $\tilde\varepsilon^m_q>0$ such that $\inf_{S\in\mathcal{U}_m(\omega^d_p)}\varepsilon_q(S)\ge \tilde\varepsilon^m_q$.
\begin{proof}
For fixed $m$, by the first bullet of \textbf{Step 2},
\[
K'_m:=\Bigl\{k:\ \exists\,S\in\mathcal{U}_m(\omega^d_p)\ \text{with}\ \overline{\mathbf{B}}^{\mathbf F}_{\lambda+1,\{k\}}\cap \mathbf{B}^{\mathbf F}_{\eta(S)}(S)\neq\emptyset\Bigr\}
\subset \bigcup_{k=1}^{k_m}\bar K(S_k),
\]
hence $\#K'_m<\infty$ by the third bullet of \textbf{Step 1}. By the inductive hypothesis and the second bullet of \textbf{Step 2}, the ratios $\eta(S)/a_q(S)$ admit a positive lower bound $c'_m>0$ on $\mathcal{U}_m(\omega^d_p)$. Therefore,
\[
\begin{aligned}
\inf_{S\in\mathcal{U}_m(\omega^d_p)}\varepsilon_q(S)
&= \min_{k\in K'_m}\ 
\inf_{\substack{S\in\mathcal{U}_m(\omega^d_p)\\
\overline{\mathbf{B}}^{\mathbf F}_{\lambda+1,\{k\}}
\cap \mathbf{B}^{\mathbf F}_{\eta(S)}(S)\neq\emptyset}}
\frac{c_{0,k}}{16}\left[h_{S_k}\!\left(\frac{\eta(S)}{4a_q(S)}\right)\right]^2 \\
&\ge \min_{k\in K'_m}\frac{c_{0,k}}{16}\,h_{S_k}(c'_m)^2 \;>\; 0.
\end{aligned}
\]
\end{proof}

Choose $a_{q+1}(S)\ge 2a_q(S)+1$ minimal such that for every $V\in\overline{\mathbf{B}}^{\mathbf F}_{\eta(S)/a_{q+1}(S)}(S)$,
\[
\|V\|(M)\ \ge\ \omega^d_p-\varepsilon_q(S).
\]

\textbf{Claim 3.} There exists $\tilde a^m_{q+1}>0$ such that $\sup_{S\in\mathcal{U}_m(\omega^d_p)} a_{q+1}(S)\le \tilde a^m_{q+1}$.
\begin{proof}
Since $\eta\le 1$, it suffices to find $c''_m>0$ such that for every $S\in\mathcal{U}_m(\omega^d_p)$ and every $V\in\overline{\mathbf{B}}^{\mathbf F}_{c''_m}(S)$,
\[
\|V\|(M)\ \ge\ \omega^d_p-\tilde\varepsilon^m_q.
\]
Then $c''_m\le \eta(S)/a_{q+1}(S)\le 1/a_{q+1}(S)$, so $a_{q+1}(S)\le 1/c''_m$ uniformly on $\mathcal{U}_m(\omega^d_p)$.

Argue by contradiction: suppose there exist $S_i\in\mathcal{U}_m(\omega^d_p)$ and $V_i\in\overline{\mathbf{B}}^{\mathbf F}_{1/i}(S_i)$ with $\|V_i\|(M)<\omega^d_p-\tilde\varepsilon^m_q$. By compactness of $\mathcal{U}_m(\omega^d_p)$, up to a subsequence $S_i\to S$ and $V_i\to S$, so $\|S\|(M)=\omega^d_p$, a contradiction.
\end{proof}

\textbf{Claim 4.} If $\Theta(x)\notin \mathbf{B}^{\mathbf F}_{\eta(S)/a_q(S)}(S)$ and $\|\Theta(x)\|(M)-\omega^d_p\le \varepsilon_q(S)$, then
\[
H^{\Theta,\delta}_{\lambda,K}(x,t)\notin \mathbf{B}^{\mathbf F}_{\eta(S)/a_{q+1}(S)}(S).
\]
\begin{proof}
For $t\in[0,1/2]$, since $\delta<\eta(S)/(4a_q(S))$ and by the third bullet of \textbf{Step 2},
\[
\begin{aligned}
\mathbf{F}\bigl(H^{\Theta,\delta}_{\lambda,K}(x,t),S\bigr)
&\ge \mathbf{F}\bigl(\Theta(x),S\bigr)
 - \mathbf{F}\bigl(H^{\Theta,\delta}_{\lambda,K}(x,t),\Theta(x)\bigr) \\
&\ge \frac{\eta(S)}{a_q(S)} - \delta
\ \ge\ \frac{3\,\eta(S)}{4\,a_q(S)}
\ \ge\ \frac{\eta(S)}{a_{q+1}(S)}.
\end{aligned}
\]


For $t\in(1/2,1]$, we have $\delta<\varepsilon_q(S)$, hence by the third bullet of \textbf{Step 2},
$\|H^{\Theta,\delta}_{\lambda,K}(x,1/2)\|(M)\le \omega^d_p+2\varepsilon_q(S)$. If, toward a contradiction, there exists $t_1\in(1/2,1]$ with
$H^{\Theta,\delta}_{\lambda,K}(x,t_1)\in\overline{\mathbf{B}}^{\mathbf F}_{\eta(S)/a_{q+1}(S)}(S)$, let $k=\bar k(K)$. Then
\[
\begin{aligned}
\mathrm{dist}\!\Big(
\bigl((F^k_{\cdot})_{\#}(\Theta(x))\bigr)^{-1}\!\bigl(&H^{\Theta,\delta}_{\lambda,K}(x,t_1)\bigr),\ 
\bigl((F^k_{\cdot})_{\#}(\Theta(x))\bigr)^{-1}\!\bigl(H^{\Theta,\delta}_{\lambda,K}(x,1/2)\bigr)
\Big) \\
&\ge d_{q,S,k}(\Theta(x))
\ \ge\ h_{S_k}\!\left(\frac{\eta(S)}{4a_q(S)}\right) \;>\; 0.
\end{aligned}
\]

By Lemma~\ref{lem:grad}, in addition to equations \eqref{eq:flow} and \eqref{eq:def-eps} above,
\[
\|H^{\Theta,\delta}_{\lambda,K}(x,t_1)\|(M)
\ \le\ \|H^{\Theta,\delta}_{\lambda,K}(x,1/2)\|(M)-\frac{c_{0,k}}{2}\left[h_{S_k}\!\left(\frac{\eta(S)}{4a_q(S)}\right)\right]^2
\ \le\ \omega^d_p-6\varepsilon_q(S),
\]
contradicting the defining property of $a_{q+1}(S)$.
\end{proof}

Finally, define
\[
\varepsilon_m:=\min_{q\in\{1,\dots,p(n-d)\}}\ \inf_{S\in\mathcal{U}_m(\omega^d_p)}
\min\!\left(\frac{\eta(S)}{4a_{q+1}(S)},\,\varepsilon_q(S)\right)>0.
\]

\textbf{Step 4 [Hierarchical Deformations].}
In this step we construct a (sub)sequence of homotopies
\[
H_i: X_i \times [0,1] \longrightarrow \mathcal Z_d(M; \mathbb Z_2)
\]
such that \(H_i(\cdot,0)=\Phi_i\) and \(H_i(\cdot,1)=\Psi_i\), with all required properties.

Choose a subsequence \(\{j_i\}\) so that \(\sup_x \mathbf M(\Phi_{j_i}(x)) \le \omega^d_p + \varepsilon_i/2\), which is possible by the definition of a min--max sequence. Relabel this subsequence as \(\Phi_i\).

It suffices to deform \(\Phi_i\) to \(\Psi_i\) so that
\(|\Psi_i|\cap \mathcal N^i_{\eta/(2a_{p(n-d)+1})}=\emptyset\), where for
\(c:\mathcal U(\omega^d_p)\to\mathbb R^+\) we set
\[
\mathcal N^i_{c\eta}\ :=\ \bigcup_{S\in\mathcal U_i(\omega^d_p)} \mathbf B^{\mathbf F}_{\,2c(S)\eta(S)}(S).
\]
Fix \(i\) for the remainder of this step.

Recall that subdivision of a simplicial complex does not change the topological or metric properties, only the combinatorial structure. We perform repeated barycentric subdivisions of \(X_i\) to obtain a “nice’’ combinatorial structure and partition the closed faces into \textit{good} and \textit{bad} faces, with the following properties:
\begin{enumerate}
  \item If \(|\Phi_i|(F)\cap \mathcal N^i_{\eta/2}\neq\emptyset\), then \(|\Phi_i|(F)\subset \mathcal N^i_\eta\). Denote by \(\mathcal B\) the set of all such \textit{bad} faces.
  \item For each \(F\in\mathcal B\) there is a nonempty finite set \(K(F)\subset\mathbb N\) such that \(|\Phi_i|(F)\subset \mathbf B^{\mathbf F}_{K(F)}\), and whenever \(F\subset F'\) with \(F,F'\in\mathcal B\) we have \(K(F)\supset K(F')\).
\end{enumerate}
Since \(X_i\) is compact, \(|\Phi_i|\) is uniformly continuous. Moreover, \(\eta\) has a uniform positive lower bound on \(\mathcal U_i(\omega^d_p)\); hence \(\mathbf F(\partial\mathcal N^i_\eta,\mathcal N^i_{\eta/2})\ge c_i>0\). After sufficiently many barycentric subdivisions, property (1) holds.

By the first bullet of \textbf{Step 2}, each $|\Phi_i|(F)$ with $F\in\mathcal B$ can be covered by finitely many
\begin{equation*}
    \{\mathbf{B}^\mathbf{F}_{\varepsilon_j}(S_j)\}.
\end{equation*}
Hence the pullback open cover
\begin{equation*}
    \{\,|\Phi_i|^{-1}(\mathbf{B}^\mathbf{F}_{\varepsilon_j}(S_j))\,\}
\end{equation*}
is finite and therefore admits a Lebesgue number. Because $\mathcal B$ is finite, after further barycentric subdivisions each $|\Phi_i|(F)$ with $F\in\mathcal B$ lies in one of the $\mathbf{B}^\mathbf{F}_{\varepsilon_j}(S_j)$.

\begin{figure}
\centering

\caption{Obtaining a homotopy on $F'_k$ from one on $Y$}
\label{YFkk}
\end{figure}

We define the index sets $K$ as follows. If $F\in\mathcal B$ is not contained in any other member of $\mathcal B$, set $K(F)$ to be the singleton consisting of one index $j$ such that $\mathbf{B}^\mathbf{F}_{\varepsilon_j}(S_j)$ contains $|\Phi_i|(F)$. Proceeding by downward induction on dimension, for each face $F'\subset F$ with $F\in\mathcal B$, define
\begin{equation*}
K(F') := \bigcup_{F'\subset F\in\mathcal B} K(F).
\end{equation*}
By construction, this has the required properties. We cannot simply set $K(F')=K(F)$ because a lower-dimensional face $F'$ may lie in multiple higher-dimensional faces $F$ in $\mathcal B$, which can yield different sets.

We now define \(H_i\) on skeleta \(X_i^{(k)}\) by induction on \(k\). For $k = 0$, we apply \textbf{Step 2} with $\lambda = 0$ and $\delta < \varepsilon_i/4$ to all the $0$-cells in $\mathcal{B}$. For all the $0$-cells outside $\mathcal{B}$, we simply use the constant homotopy map. This yields $H_i^{(0)}: X_i^{(0)}\times[0,1]\to \mathcal Z_d(M;\mathbb Z_2)$ such that
\begin{equation*}
    |H_i^{(0)}|(X_i^{(0)}\times\{1\})\ \cap\ \mathcal N^i_{\eta/a_1}\ =\ \emptyset,
    \qquad
    |H_i^{(0)}|(x,[0,1])\ \subset\ \mathbf B^{\mathbf F}_{\,1,K(x)}
    \ \text{ for }x\in\mathcal B^{(0)}.
\end{equation*}

We assume the homotopy map has been defined on the $(k-1)$-skeleton, with $k\ge 1$. Consider a $k$-face $F_k$. Set
\begin{equation*}
F'_k := F_k \cup \bigl(\partial F_k \times [0,1]\bigr).
\end{equation*}
Then $F'_k$ is homeomorphic to $F_k$ ($\cong D^k$), so we can define $\Theta$ on $F'_k$ by concatenating $H^{(k-1)}_i$ on $\partial F_k \times [0,1]$ (available by the inductive hypothesis) with $\Phi_i$ on $F_k$. Now we construct a homotopy map $\tilde H^{(k)}_i$ with initial data $\Theta$. We have two cases:

\textbf{Case 1.} If $F_k \notin \mathcal{B}$, then according to the definition of $\mathcal{B}$, no cell in $\partial F_k$ belongs to $\mathcal{B}$. In this case we define the homotopy on this cell to be constant.

\textbf{Case 2.} If \(F_k\in\mathcal B\), the subdivision guarantees
\(|\Theta|(F_k')\subset \mathbf B^{\mathbf F}_{\,k,K(F_k)}\).

However, it may fail that \(|\Theta|(F_k')\subset\mathcal N^i_\eta\). Subdivide \(F_k'\) further so that every closed \(k\)-face \(\widetilde F\) with \(|\Theta|(\widetilde F)\cap \mathcal N^i_{\eta/a_k}\neq\emptyset\) actually satisfies \(|\Theta|(\widetilde F)\subset \mathcal N^i_\eta\). Let \(Y\) be the union of all such \(\widetilde F\). By induction, \(|\Theta|(\partial F_k')\cap \mathcal N^i_{\eta/a_k}=\emptyset\), hence \(|\Theta|(\partial Y)\cap \mathcal N^i_{\eta/a_k}=\emptyset\).

With this new subdivision, we can define a map 
\begin{equation*}
    \hat H: F'_k \times [0, 1] \cup Y \times [1,2] \rightarrow \mathcal{Z}_d(M; \mathbb{Z}_2) 
\end{equation*}
\begin{figure}
\centering

\caption{Obtaining a homotopy on $F_k$ from one on $F'_k$}
\label{F'Fk}
\end{figure}
such that $\hat H(\cdot, t) := \Theta(\cdot)$ for $t \in [0,1]$, and $\hat H(\cdot, t+1) = H^{|\Theta|, \delta}_{k, K(F_k)}(\cdot, t)$ for $t \in [0,1]$ in \textbf{Step 2} with $\delta < \varepsilon_i$. By the construction of $Y$, we have
\begin{equation*}
    |\hat H((F'_k - Y) \times 1 \cup \partial F'_k \times [0, 1])| \cap \mathcal{N}^i_{\eta/a_{k+1}} = \emptyset\,.
\end{equation*}
By (5) in the third bullet of \textbf{Step 2},
\begin{equation*}
    |\hat H(Y \times 2)| \cap \mathcal{N}^i_{\eta/a_{k+1}} = \emptyset\,.
\end{equation*}
By \textbf{Step 3},
\begin{equation*}
    |\hat H(\partial Y \times [1, 2])| \cap \mathcal{N}^i_{\eta/a_{k+1}} = \emptyset\,.
\end{equation*} 
We can derive $\tilde H^{(k)}_i: F'_k \times [0, 1] \rightarrow \mathcal{Z}_d(M; \mathbb{Z}_2)$ from $\hat H$ induced by the homeomorphism (see Figure \ref{YFkk}.)
\begin{equation*}
    \begin{aligned}
        (F'_k \times [0,1], &\partial F'_k \times [0,1], F'_k \times 1) \cong 
        (F'_k \times [0, 1] \cup Y \times [1,2], \\
        & \partial F'_k \times [0,1], (F'_k\backslash Y)\times 1) \cup \left(Y \times 2\right) \cup \left(\partial Y \times [1,2]\right))\,.
    \end{aligned}            
\end{equation*}
        
By the property of $\hat H$, we also have
\begin{equation*}
    |\tilde H^{(k)}_i(F'_k \times 1 \cup \partial F'_k \times [0, 1])| \cap \mathcal{N}^i_{\eta / a_{k + 1}} = \emptyset\,.
\end{equation*}
For both cases, we can derive $H^{(k)}_i: F_k \times [0, 1] \rightarrow \mathcal{Z}_d(M; \mathbb{Z}_2)$ from $\tilde{H}^{(k)}_i$ induced by the homeomorphism 
\begin{equation*}
    (F_k \times [0,1], F_k \times 1) \cong (F'_k \times [0,1], F'_k \times 1 \cup \partial F'_k \times [0,1]),
\end{equation*}
satisfying that $H^{(k)}_i|_{\partial F_k \times [0,1]} = H^{(k-1)}_i|_{\partial F_k \times [0,1]}$ and $H^{(k)}_i|_{F_k}(\cdot, 0) = \Phi_i|_{F_k}$ (see Figure \ref{F'Fk}). Concatenating over all \(k\)-cells yields
\begin{equation*}
    H^{(k)}_i: X^{(k)}_i \times [0,1] \rightarrow \mathcal{Z}_d(M; \mathbb{Z}_2).
\end{equation*}
Note that
\begin{equation*}
    |H^{(k)}_i(F_k \times 1)| = |\tilde H^{(k)}_i(F'_k \times 1 \cup \partial F'_k \times [0, 1])|\,.   
\end{equation*}
Therefore, we can conclude that
\begin{equation*}
    |H^{(k)}_i|(X^{(k)}_i \times 1) \cap \mathcal{N}^i_{\eta / a_{k+1}} = \emptyset\,.
\end{equation*}

Set \(\Psi_i:=H_i^{(p(n-d))}(\cdot,1)\) and \(\bar\varepsilon(S):=\eta(S)/(2a_{p(n-d)+1}(S))\). Then all required conclusions follow.
\end{proof}

\end{document}
